\chapter{Лабораторная работа №4. Основы сетевой безопасности и доступа к сети.}
\label{ch:chapter 4}

\section{\textit{Лабораторная работа №4.1. Настройка ACL.}}

\subsection{Цели}

Лабораторная работа помогает получить практические навыки по изучению
следующих тем:

\begin{itemize}

\item Настройка списков ACL

\item Применение ACL на интерфейсе

\item Основные методы фильтрации трафика

\end{itemize}

\subsection{Топология сети}

В сети, показанной на схеме, маршрутизатор R3 выполняет функции сервера,
маршрутизатор R1 выполняет функции клиента, и они доступны для связи.
Физические интерфейсы, соединяющие R1 и R2, имеют IP-адреса 10.1.2.1/24 и
10.1.2.2/24 соответственно, а физические интерфейсы, соединяющие R2 и R3, —
IP-адреса 10.1.3.2/24 и 10.1.3.1/24. Кроме того, на маршрутизаторе R1 созданы два
логических интерфейса LoopBack 0 и LoopBack 1 для имитации двух пользователей-
клиентов. Два интерфейса имеют IP-адреса 10.1.1.1/24 и 10.1.4.1/24 соответственно.
Один пользователь (LoopBack 1 на R1) должен удаленно управлять R3. Для гарантии
того, что вход в R3 будет разрешен только пользователю, который соответствует
политике безопасности, можно настроить Telnet на сервере, задать защиту паролем и
сконфигурировать ACL.

\includegraphics[width=0.5\textwidth]{lb4_1.png} 
\captionof{figure}{Топология сети для лабораторной работы №4}

\subsection{План работы}

\begin{itemize}

\item 1. Настройка IP-адресов.

\item 2. Настройка OSPF для обеспечения возможности сетевого подключения.

\item 3. Создание ACL на основе необходимого трафика.

\item 4. Настройка фильтрации трафика.

\end{itemize}

\subsection{Процедура конфигурирования}

\textbf{Шаг 1.} Настройка IP-адресов.

Настроим IP-адреса для R1, R2 и R3:

\includegraphics[width=0.5\textwidth]{lb4_2.png} 

\includegraphics[width=0.5\textwidth]{lb4_3.png} 

\includegraphics[width=0.5\textwidth]{lb4_4.png} 

\textbf{Шаг 2.} Настроим OSPF для обеспечения возможности сетевого подключения.

Настроим OSPF на маршрутизаторах R1, R2 и R3 и назначьте их в область 0, чтобы
обеспечить возможность подключения:

\includegraphics[width=0.5\textwidth]{lb4_5.png} 

\includegraphics[width=0.5\textwidth]{lb4_6.png} 

\includegraphics[width=0.5\textwidth]{lb4_7.png} 

Выполним команду ping на маршрутизаторе R3, чтобы проверить возможность
подключения к сети:

\includegraphics[width=0.5\textwidth]{lb4_8.png} 

\textbf{Шаг 3.} Сконфигурируем R3 в качестве сервера.

Включим функцию Telnet на R3, установите для уровня пользователя значение 3 и
задайте для входа пароль — Huawei@123:

\includegraphics[width=0.5\textwidth]{lb4_9.png} 

Команда telnet server enable позволяет включить службу Telnet.

\includegraphics[width=0.5\textwidth]{lb4_10.png} 

Команда user-interface позволяет перейти в режим интерфейса одного или
нескольких пользователей.

Пользовательский интерфейс терминала виртуального типа (Virtual Type Terminal,
VTY) осуществляет управление и мониторинг входа пользователей в систему с
помощью Telnet или SSH.

\includegraphics[width=0.5\textwidth]{lb4_11.png} 

В данном случае нужно вводить сразу после cipher, из-за устройства роутера.

\textbf{Шаг 4.} Настроим ACL на основе необходимого трафика.

Способ 1. Настроим ACL на интерфейсе VTY маршрутизатора R3, чтобы разрешить
вход с R1 в R3 через Telnet, используя IP-адрес LoopBack 1.

Настроим ACL на R3:

\includegraphics[width=0.5\textwidth]{lb4_12.png} 

Выполним фильтрацию трафика на интерфейсе VTY маршрутизатора R3:


\includegraphics[width=0.5\textwidth]{lb4_13.png} 

Выведем на экран конфигурацию ACL на R3:

\includegraphics[width=0.5\textwidth]{lb4_14.png} 

Команда display acl позволяет вывести на экран конфигурацию ACL.

Создан расширенный ACL. Он имеет номер 3000 и содержит два правила.
Правила ACL пронумерованы с шагом 5.
Правило 5 разрешает прохождение соответствующего трафика. Если пакетов,
соответствующих правилу, нет, поле matches не отображается.

Способ 2. Настроим ACL на физическом интерфейсе маршрутизатора R2, чтобы
разрешить вход с R1 в R3 через Telnet, используя IP-адрес физического интерфейса.

Настроим ACL на R2:

\includegraphics[width=0.5\textwidth]{lb4_15.png} 

Выполним фильтрацию трафика на интерфейсе GE0/0/3 маршрутизатора R3:

К сожалению, в ходе использования выявилась проблема, связанная с отсутствием необходимой команды на  роутере, вследствии чего нельзя завершить лабораторную работу.

\subsection{Вопросы.} 
R3 выполняет функции как сервера Telnet, так и FTP-сервера, IP-адрес LoopBack 0 на
R1 должен использоваться только для доступа к службе FTP, а IP-адрес LoopBack 1 на
R1 должен использоваться для удаленного управления R3 через Telnet.
Настройте ACL в соответствии с приведенными требованиями.

Вследствии вышеописанной проблемы и это невозможно сделать.

\section{\textit{Лабораторная работа №4.2. Настройка локального механизма AAA.}}

\subsection{Цели}

Лабораторная работа помогает получить практические навыки по изучению
следующих тем:

\begin{itemize}

\item Настройка локального механизма AAA

\item Процедура создания домена

\item Процедура создания локального пользователя

\item Управление пользователями на основе домена

\end{itemize}

\subsection{Топология сети}

Маршрутизатор R1 выполняет функции клиента, а R2 — функции сетевого устройства.
Необходим контроль доступа к ресурсам на R2. Следовательно, нужно настроить
локальную аутентификацию AAA на маршрутизаторах R1 и R2 и управлять пользователями на основе доменов, а также настроить уровни полномочий для
аутентифицированных пользователей.

\includegraphics[width=0.5\textwidth]{lb4_25.png} 
\captionof{figure}{Топология сети для лабораторной работы №4}

\subsection{План работы}

\begin{itemize}

\item 1. Настройка IP-адресов.

\item 2. Настройка OSPF для обеспечения возможности сетевого подключения.

\item 3. Создание ACL на основе необходимого трафика.

\item 4. Настройка фильтрации трафика.

\end{itemize}

\subsection{Процедура конфигурирования}

\textbf{Шаг 1.} Настроим IP-адреса для маршрутизаторов R1 и R2.

\includegraphics[width=0.5\textwidth]{lb4_16.png} 

\includegraphics[width=0.5\textwidth]{lb4_17.png} 

\textbf{Шаг 2.} Настроим схему AAA.

Настроим схемы аутентификации и авторизации:

\includegraphics[width=0.5\textwidth]{lb4_18.png} 

Устройство, выполняющее функции сервера AAA, называется локальным сервером
AAA, который может выполнять аутентификацию и авторизацию, но не учет.

Локальному серверу AAA требуется база данных локальных пользователей,
содержащая имя пользователя, пароль и информацию об авторизации локальных
пользователей. Локальный сервер AAA быстрее и дешевле, чем удаленный сервер
AAA, но имеет меньшую емкость хранилища.

\textbf{Шаг 3.} Создайте домен и примените к нему схему AAA.

\includegraphics[width=0.5\textwidth]{lb4_19.png} 

Устройства управляют пользователями на основе доменов. Домен — это группа
пользователей. Каждый пользователь принадлежит к домену. Конфигурация AAA для
домена применяется к пользователям в домене. Создайте домен с именем datacom.

\textbf{Шаг 4.} Настроим локальных пользователей.

Создадим локального пользователя и настройте для него пароль:

\includegraphics[width=0.5\textwidth]{lb4_20.png} 

Если в имени пользователя содержится разделитель в виде символа @, то строка
символов перед символом @ — это имя пользователя, а строка символов после
символа @ — имя домена. Если в имени пользователя нет разделителя в виде
символа @, вся символьная строка представляет собой имя пользователя, а для
домена используется имя по умолчанию.

Настроим параметры для локального пользователя, такие как тип доступа и
уровень полномочий:

\includegraphics[width=0.5\textwidth]{lb4_21.png} 


Команда local-user service-type позволяет настроить тип доступа для локального
пользователя. После настройки типа доступа пользователь сможет успешно войти в
систему, только применив настроенный тип доступа. Если в качестве типа доступа
было указано telnet, пользователь не сможет получить доступ к устройству через веб-
страницу. Для одного пользователя можно настроить несколько типов доступа.

Для локального пользователя настроен уровень полномочий. Этому пользователю
будут доступны только команды, предоставляемые указанным уровнем полномочий,
и команды более низкого уровня полномочий.

\textbf{Шаг 5.} Включим функцию telnet на R2.

\includegraphics[width=0.5\textwidth]{lb4_22.png}

Команда authentication-mode позволяет настроить режим аутентификации для
доступа к пользовательскому интерфейсу. По умолчанию режим аутентификации на
пользовательском интерфейсе VTY не настроен. Для интерфейса входа в систему
необходимо настроить режим аутентификации. В противном случае пользователи не
смогут выполнять вход в устройство.

\textbf{Шаг 6.} Проверим конфигурацию AAA.

Выполним вход с R1 на R2 через Telnet:

\includegraphics[width=0.5\textwidth]{lb4_23.png}

Выведем на экран список пользователей, подключенных к R2:

\includegraphics[width=0.5\textwidth]{lb4_24.png}



\section{\textit{Лабораторная работа №4.3. Настройка NAT.}}

\subsection{Цели}

Лабораторная работа помогает получить практические навыки по изучению
следующих тем:

\begin{itemize}

\item Настройка динамического NAT

\item Настройка Easy IP

\item Настройка NAT-сервера

\end{itemize}

\subsection{Топология сети}

Для решения проблемы нехватки адресов IPv4 предприятия, как правило, используют
частные адреса IPv4. Однако корпоративная сеть должна предоставлять доступ
сотрудникам к общедоступной сети и услуги внешним пользователям. В этом случае
необходимо настроить NAT в соответствии с приведенными выше требованиями.

1. Сеть между маршрутизаторами R1 и R2 является интрасетью и использует
частные адреса IPv4.

2. R1 выполняет функции клиента, а R2 является шлюзом для R1 и граничным
маршрутизатором, подключенным к общедоступной сети.

3. R3 имитирует общедоступную сеть.

\includegraphics[width=0.5\textwidth]{lb4_26.png} 
\captionof{figure}{Топология сети для лабораторной работы №4}

\subsection{План работы}

\begin{itemize}

\item 1. Настройка динамического NAT.

\item 2. Настройка Easy IP.

\item 3. Настройка сервера NAT.

\end{itemize}

\subsection{Процедура конфигурирования}

\textbf{Шаг 1.} Настроим основные параметры.

Настроим IP-адреса и маршруты:

\includegraphics[width=0.5\textwidth]{lb4_27.png} 

\includegraphics[width=0.5\textwidth]{lb4_28.png} 

\includegraphics[width=0.5\textwidth]{lb4_29.png} 

Настроим функцию Telnet на маршрутизаторах R1 и R3 для последующей проверки:

\includegraphics[width=0.5\textwidth]{lb4_30.png} 

\includegraphics[width=0.5\textwidth]{lb4_31.png} 

Проверим возможность установления связи:

\includegraphics[width=0.5\textwidth]{lb4_32.png} 

\includegraphics[width=0.5\textwidth]{lb4_33.png} 

У маршрутизатора R1 нет связи с R3, потому что на R3 не настроен маршрут к адресу
192.168.1.0/24.

Более того, на R3 нельзя настраивать маршруты в частные сети.

\textbf{Шаг 2.} Предприятие получает общедоступные IP-адреса в диапазоне от 1.2.3.10 до 1.2.3.20,
поэтому ему требуется функция динамического NAT.

Настроим пул адресов NAT:

\includegraphics[width=0.5\textwidth]{lb4_34.png} 

С помощью команды nat address-group можно настроить пул адресов NAT. В данном
примере пул адресов имеет номер 1 Пул адресов должен быть набором
последовательных IP-адресов. При достижении внутренними пакетами данных
границы частной сети частные IP-адреса источников будут преобразовываться в
общедоступные IP-адреса.

Настроим ACL:

\includegraphics[width=0.5\textwidth]{lb4_35.png} 

Настроим динамический NAT на GigabitEthernet0/0/4 маршрутизатора R2:

\includegraphics[width=0.5\textwidth]{lb4_36.png} 

Команда nat outbound позволяет установить привязку ACL к пулу адресов NAT.
IP-адреса пакетов, соответствующих списку ACL, будут преобразовываться в адреса
из пула адресов. Если в пуле достаточно адресов, можно добавить аргумент no-pat,
чтобы включить однозначное преобразование адресов. В этом случае будут
преобразовываться только IP-адреса пакетов данных, а порты преобразовываться не
будут.

Выведем на экран таблицу сеансов NAT на R2:

\includegraphics[width=0.5\textwidth]{lb4_37.png} 

\textbf{Шаг 3.} Если IP-адрес GigabitEthernet0/0/4 на R2 назначается динамически (например, через
DHCP или PPPoE), необходимо настроить Easy IP.

Удалим конфигурацию, созданную на предыдущем шаге:

\includegraphics[width=0.5\textwidth]{lb4_38.png} 

Настроим Easy IP:

\includegraphics[width=0.5\textwidth]{lb4_39.png} 

Проверим возможность установления связи:

\includegraphics[width=0.5\textwidth]{lb4_40.png} 

\textbf{Шаг 4.} R3 должен предоставлять сетевые услуги (в данном примере telnet) для
пользователей в общедоступной сети. Поскольку R3 не имеет общедоступного
IP-адреса, необходимо настроить сервер NAT на исходящем интерфейсе R2.

Настроим сервер NAT на R2:

\includegraphics[width=0.5\textwidth]{lb4_41.png} 

Команда nat server позволяет определить таблицу сопоставления внутренних
серверов, чтобы внешние пользователи могли получать доступ к внутренним серверам через преобразование адресов и портов. Можно настроить внутренний
сервер так, чтобы пользователи внешней сети могли инициировать доступ к
внутреннему серверу. Когда хост во внешней сети отправляет запрос на соединение
на общедоступный адрес (глобальный адрес) внутреннего сервера NAT, сервер NAT
преобразует адрес назначения, содержащийся в запросе, в частный адрес
(внутренний адрес) и пересылает запрос на сервер в частной сети.

Выполним вход с R3 на R1 через Telnet:

\includegraphics[width=0.5\textwidth]{lb4_42.png} 

\includegraphics[width=0.5\textwidth]{lb4_43.png} 

\textbf{Вопросы.}1. После настройки сервера NAT должны ли порты назначения до преобразования
соответствовать портам назначения после преобразования?

-Нет, порты назначения до и после преобразования не обязательно должны совпадать.



